\section{The measure, the record, and the conventions}\label{sec:construct}\label{sec:measure}

\subsection{The record and the question}

A US trademark application carries a description of the goods or services
the mark will cover. The description is written to secure scope of
protection: wide enough to cover the planned use, narrow enough to be
granted. It is therefore a forced, dated disclosure of what a firm expected
to sell. The USPTO publishes every application since 1884 with its
prosecution history, and this paper re-parses all 13.99 million case files
into dated events, so that each application can be followed from filing to
registration, to the sworn proof of continued use required between the fifth
and sixth year, and to cancellation or renewal.

The question is whether the language of the description, read against the
language of the rest of the industry at the time, predicts what happened to
the mark and to the firm behind it. Two properties of the language are
measured. \emph{Atypicality} is how unusual the description is for its
industry. \emph{Lead} is whether its language is ahead of or behind where the
industry's language was moving. Sections~\ref{sec:registration}--\ref{sec:success}
relate the two to reaching registration, to the five-year proof of continued
use, and to venture funding and public listing. Section~\ref{sec:robust}
reports what changes under other choices.

\subsection{Two readers, one living backwards}

Imagine two readers of a filing's goods/services description, each having seen
a disjoint half of its industry. One has read only what the industry filed in
the five years before it. The other, living backwards as Merlin was said to,
has read only the five years after. Each is asked how surprising the wording
is. The measurement is how much they disagree: a description that startles the
first reader and not the second arrived before its industry's language moved
that way.

Both windows are anchored on the filing's own date, spanning the 1,826
days before it
and the 1,826 days after; both exclude the filing date itself, so a filing
never appears in the reference of another filing made the same day, and no
two filings made on different days share a reference. Inside a window every day counts the same. The measure asks whether
wording is unusual against a period, not whether it is ahead of a trend
inside that period; it has no momentum channel.

\subsection{What is measured}

The quantity is information content in Shannon's sense. Something that
happens with probability $p$ carries $-\log p$ of information when it
happens: a common thing carries little, a rare thing a lot. Here the things
are the themes a description draws on, and the probabilities are how often
the filing's own industry drew on each of them in the reference window. A
filing's score is the average information content of its themes under its
industry's frequencies, weighted by how heavily it uses each, less the same
average taken under the filing's own frequencies.

The apparatus is that of \citet{Murdock2017}, developed on Darwin's reading
notebooks, and \citet{Barron2018}, who extended it to French Revolution
parliamentary debates. They score a document against those before it and
again against those after, calling the two quantities \emph{novelty} and
\emph{transience} and their signed difference \emph{resonance}. Their object
was the life of ideas: how often an idea appeared, and how ideas stood in
relation to one another, across one reading life or one assembly's debate.
The object here is the composition of commercial offerings, which components
a product or service is described as having, and how that composition shifts
across an industry over time. The arithmetic carries over unchanged. The
interpretation does not: their document has one author and the reference
stream is that author's own reading, so a high value says something about a
mind; here a filing is scored against other firms' filings, and its author is
usually counsel drafting for scope of protection.

\paragraph{Themes, not topics.} The clusters the model fits are called
\emph{themes} throughout. In the source literature and in the name of the
method, latent Dirichlet allocation, the same objects are called topics, and
a document's topic is what it is about. A goods/services description is not
about anything in that sense; it lists components. A cluster of terms that
co-occur across descriptions (\emph{video, music, games, electronic,
audio}) is a component that many offerings share, not a subject. Theme is
the word for that. Topic is used here only in the name of the model.

\subsection{From a description to two numbers}\label{sec:definitions}

\paragraph{Counting words.} A description is reduced to a bag of words: the
text is lowercased, split into tokens of three or more letters, and counted,
with adjacent pairs counted as tokens in their own right so that
\emph{computer software} registers alongside \emph{computer} and
\emph{software}. Order is discarded. Take the filing for \textsc{Amazon.com},
lodged in the advertising and retail class on 18 April 1997:

\begin{quote}\small
computerized on line search and ordering service featuring the wholesale and
retail distribution of books, music, motion pictures, multimedia products and
computer software in the form of printed books, audiocassettes,
videocassettes, compact disks, floppy disks, CD ROMs, and direct digital
transmission
\end{quote}

\noindent Of its 45 distinct terms, 36 are in the theme model's vocabulary,
among them \emph{computerized}, \emph{search}, \emph{ordering},
\emph{ordering service}, \emph{retail}, \emph{distribution}, \emph{printed
books} and \emph{digital transmission}.

\paragraph{Grouping words into themes.} Each filing is then re-described in a
fixed, much smaller set of coordinates. Latent Dirichlet allocation is run
once over a stratified sample of the whole corpus, all 45 classes together: it
learns 50 clusters of words that tend to occur together and, for any filing,
the proportions in which its words draw on them. Nothing supervises this; the
clusters come from co-occurrence alone. This build --- fifty global themes,
per-filing five-year windows, flat weighting --- is the production scoring,
used for every estimate unless the text states another. The themes are
shared across classes;
what is specific to a class is the reference, the average theme proportions
of that class's own filings in the window. The Amazon filing draws on six
themes at more than a twentieth of its weight: $0.41$ on media and
entertainment goods (\emph{video, computer, electronic, music, games,
audio}), $0.15$ on computer and information services (\emph{services,
computer, software, information, data}), $0.13$ on retail store services
(\emph{store, retail, store services, featuring}), $0.10$ on printed matter
(\emph{paper, books, printed, cards, stationery}), $0.09$ on downloadable and
online software, and $0.07$ on marketing of consumer goods. The online
appendix lists every theme with its leading words and its weight in each
class.\footnote{\url{https://aporia.institute/tm-vocabulary/online-appendix/themes_T50.html}}

Fifty is coarse, roughly one theme per industry.
Section~\ref{sec:three_scorings} refits the model at 500 themes and
separately at 50 themes within each class and repeats the five-year-proof
analysis under each; Section~\ref{sec:robust} reports the resolution sweep
and a second fit at the same resolution with a different seed. Fifty is
kept elsewhere because it is the coarsest resolution at which every result
holds, and because its software-class lead contrast sits at the low end of
the resolutions tried (only the 100-theme refit returns less).

\paragraph{Comparing against the two references.} Write $P$ for the filing's
distribution over themes, and $Q_{\mathrm{past}}$ and $Q_{\mathrm{future}}$
for that distribution averaged over every same-class filing in the two
windows, $[d_i - W,\, d_i)$ and $(d_i,\, d_i + W]$, with $W = 1{,}826$ days.
How far $P$ sits from a reference $Q$ is the Kullback--Leibler divergence,
\[
  \mathrm{KL}(P \,\|\, Q) \;=\; \sum_{k} P_k \log \frac{P_k}{Q_k},
\]
a sum over the 50 themes, zero when the two distributions agree and growing as
the filing puts weight where the industry puts little. It is measured in nats,
the natural-log unit. Write
\[
  K^{-} = \mathrm{KL}(P \,\|\, Q_{\mathrm{past}}),
  \qquad
  K^{+} = \mathrm{KL}(P \,\|\, Q_{\mathrm{future}}),
\]
so $K^{-}$ is what the first reader feels and $K^{+}$ the second: the
filing's past-facing and future-facing surprise.

\paragraph{Taking the difference.} Amazon's filing is scored against 41{,}166
earlier filings and 146{,}532 later ones, giving $K^{-} = 1.286$ and
$K^{+} = 1.163$. The two axes used throughout are their average and their
signed difference:
\[
  \underbrace{A \;=\; \tfrac{1}{2}\left(K^{-} + K^{+}\right)}_{\text{atypicality}},
  \qquad
  \underbrace{L \;=\; K^{-} - K^{+}}_{\text{lead}} .
\]
For Amazon, $A = 1.225$ and $L = +0.123$: against its own class, its
atypicality is at the 47th percentile and its lead at the 94th. The
industry moved toward this description in the five years after it was filed
and had not been writing that way in the five years before.
Table~\ref{tab:nomenclature} maps the terms and their nearest occupants
in the literature; the signed difference is written $L$ throughout and
its magnitude $|L|$.

What ``early'' means is visible in the themes themselves: the three themes
that carry most of Amazon's description --- retail store services, media
and entertainment goods, and the still-small online-software theme --- all
became more common in class 035 over the following five years, the last
nearly doubling, while the two minor themes it touched receded. The filing
is closer to the class's future composition than to its past one; the
supplementary material tabulates the six theme shares.

$K^{-}$ and $K^{+}$
correlate at $0.988$ on the software class, so entering both is close to
entering one variable twice; rotating them into an average and a difference
isolates the large common component from the small residual that carries the
timing, and leaves the two regressors nearly uncorrelated ($+0.036$ against
each other, where $K^{-}$ alone would give $+0.114$). The cost is that $L$ is
a small residual of two large quantities, about a sixth of either level's
spread, which Section~\ref{sec:robust} quantifies.


\paragraph{Why the words are not scored directly.} The themes are an extra
step, and the words themselves could be compared to the reference: build the
filing's distribution over the class vocabulary from the words it uses, do
the same for every filing in the window, and take the divergence. The paper
reports that scoring as a cross-check. It cannot be the primary one, because
a divergence measured on words mostly counts the length of the filing. The
divergence is an average surprise less a spread: for any two distributions,
\[
  \mathrm{KL}(P \,\|\, Q) \;=\; \underbrace{-\textstyle\sum_k P_k \log Q_k}_{\text{how unexpected these words are}}
  \;-\; \underbrace{\left(-\textstyle\sum_k P_k \log P_k\right)}_{\text{how spread out the filing is}} ,
\]
and the spread of a filing over words is bounded by how many distinct words it
has: a description using $k$ terms cannot spread over more than $k$ of them,
so its spread cannot exceed $\log k$, while the reference is spread over about
900 terms. A five-word description is charged the whole difference in width,
roughly five nats, before any of its content is considered. In 3{,}000
software-class filings of at least 150 words, scored against one fixed
reference and rescored after truncation to a random handful of words,
across a fiftyfold cut in length
the surprise term moves by $0.016$ nats and the spread term by $2.9$; synthetic filings drawn at random from the
reference itself track $\log(\text{distinct terms})$ to within a fiftieth
of a nat (supplementary material). Term-scored atypicality correlates with
the log of a filing's distinct term count at $-0.68$ in the software
class.

Themes break the coupling because the 50 coordinates are always occupied.
Under term scoring the effective number of coordinates a filing spreads over
rises seventeenfold from the shortest descriptions to the longest; under
themes it varies by less than a third, and not monotonically: a three-word
fragment gives the model little evidence, so its inferred proportions stay
close to the prior and spread across many themes. The spread term no longer
tracks length, and the score is dominated by the surprise term.
Theme-scored atypicality correlates with the log term count at $-0.22$ in
the software class and $-0.10$ in advertising and retail, against $-0.68$
and $-0.62$ under terms.

\paragraph{What the themes see, and what they do not.} The score compares a
filing's theme proportions to the average proportions of its industry, so it
reads which themes a filing draws on and how heavily. It does not, in any
material way, read how a filing combines them. Taking each filing's two
leading themes and asking how often that pair occurs together relative to the
two themes' separate frequencies, on a 150{,}000-filing software-class
sample: how rare the themes are individually explains
70\% of the variance in atypicality, and the unusualness of the pairing adds
under two points on top of it. A filing assembled from ordinary parts in a
genuinely new configuration is therefore close to invisible to this measure.

Combinations can be scored in their own right, at two levels, and the two
levels run in opposite directions (both computed on the software class,
net of the filing's own lead and atypicality; construction detail in the
supplementary material). A \emph{pair lead} --- whether the class moved
toward the filing's arrangements of themes --- is penalised like lead
itself: the top fifth fails the five-year proof of continued use $2.3$
points more often than the bottom ($t = 8.4$), and it correlates with
lead at only $-0.18$. Word-level \emph{recombination} --- adjacent word
pairs new to the class while both words are established there, novelty
assembled from familiar components \citep{Fleming2001, Uzzi2013} --- is
rewarded: the top fifth fails $2.8$ points less with lead and atypicality
held ($t = -10.4$), and the two measures are uncorrelated ($r = -0.01$).
Arranging themes the industry had not yet arranged is punished like
arriving early; coining phrases from the industry's own words is
rewarded. The construct sees neither, and the two run in opposite
directions, so its blindness to combination is a bounded omission rather
than a hidden bias in one direction.

\subsection{Within class, not across it}

Goods/services descriptions are drafted to secure scope of protection, not to
describe a product distinctively, and counsel reuse language from the USPTO's
own Acceptable Identification of Goods and Services Manual, from house
precedent, and from competitors' successful filings. Identical wording across
unrelated firms is therefore common, and it carries through to the scores:
17\% of scored registrations issued 2002--2018 share a value with another
filing in the same class and registration year. Within a class, the measure picks up departure from a shared
drafting convention. Across classes it does not: a class where the Manual
supplies dense standard language shows a compressed distribution of
atypicality, and one where firms describe novel services in prose a wide one,
for reasons unrelated to how innovative either industry is. Every estimate is
therefore taken within class and year, and magnitudes compare across classes
only as signs and orderings.

The vocabulary is positional rather than causal. A leading filing was ahead of
its industry's language, which implies a trajectory without asserting the
filing caused it. Atypicality is textual, and is not trademark law's
distinctiveness doctrine, which governs registrability along the
generic-to-fanciful spectrum; the two attach to different objects, one to the
mark and one to the description of the goods it covers.

\subsection{The gates, the samples, and the conventions}\label{sec:gates}

The outcome sections that follow are written to be read independently, so
the record's steps, the outcome names, and the counting conventions are
set out once, here.

A US trademark passes through a sequence of dated, high-stakes steps, and
each step leaves a record. An application is filed with a goods/services
description; an examiner reviews it, often issuing office actions the
applicant must answer; if allowed and (for intent-to-use filings) a sworn
statement of use is filed --- the proof that the mark has entered commerce
--- the mark registers. Registration starts a maintenance clock: in years
five to six the owner must file a Section~8 declaration proving continued
use in commerce, with a six-month grace period --- the \emph{five-year
proof of continued use}, the paper's central outcome; around year ten a
combined Section~8 declaration and Section~9 renewal is due, and every ten
years thereafter. A mark that misses a proof is cancelled. Madrid Protocol
registrations (\S66(a) basis) file Section~71 declarations on the same
schedule and are flagged separately throughout. \citet{Graham2013}
describe the administrative data; the event-dated corpus built here adds
the full prosecution history of each case, so a cancellation is assigned
to a specific proof by its date rather than inferred from a status
snapshot.\footnote{Cross-checked against the USPTO status-code dictionary:
700 registered; 701--705 Section~8 accepted; 710 cancelled for Section~8
failure; 800 renewed; 900 expired.}

\begin{table}[h]\centering\small
\caption{The trademark funnel in this corpus. Registration counts are
unique serials; five-year-proof outcomes are event-dated (a Section~8
cancellation at registration age 4.0--8.5 years). Source: the re-parsed
USPTO application backfile, scored as in this section.}
\label{tab:funnel}
\begin{tabular}{lrr}
\toprule
Stage & $n$ & Share \\
\midrule
Debut filings scored, 1995--2018 (owners' first filings) & 3{,}589{,}068 & \\
\quad reached registration & 2{,}066{,}600 & 57.6\% \\
\midrule
All unregistered class-records filed 1995--2018 & 3{,}798{,}318 & \\
\quad never filed statement of use & & 42.6\% \\
\quad stopped responding in prosecution & & 45.3\% \\
\quad removed adversarially & & 4.3\% \\
\midrule
Registrations 2002--2018, scored, unique serials & 3{,}104{,}534 & \\
\quad failed the five-year proof of continued use (event-dated) & & 52.3\% \\
\midrule
Passed the five-year proof, cohorts 2002--2013, status known & 1{,}129{,}724 & \\
\quad failed at the year-ten renewal & & 39.1\% \\
\bottomrule
\end{tabular}
\end{table}

Three named outcomes recur (Table~\ref{tab:funnel}). \emph{Failing to
register}: the application never completed, which happens by the
applicant's own inaction in roughly nine cases of ten
(Section~\ref{sec:registration}). \emph{Failing the five-year proof of
continued use}: a registration cancelled for non-use at registration age
4.0 to 8.5 years, which means the product the mark named left commerce
within six years (Sections~\ref{sec:gate} and~\ref{sec:waves}).
\emph{Reaching the SEC record}: the owner appears in dated SEC events ---
a priced Form~D round, financial reporting, a listing
(Section~\ref{sec:success}).

Outcomes are rates, so every estimate is a shift in a probability and is
stated in percentage points against the base rate it shifts: $+1.4$pp on
a 52.3\% base means 53.7\% against 52.3\%. The basic comparison is the
outcome rate of the fifth of filings with the highest lead minus that of
the fifth with the lowest, with the fifths cut inside a single Nice class
and year so that no difference between industries or between years
enters; regression tables add drafting and filer controls and cluster
errors on the owner. Where a continuous version is given, it is the
change in the outcome for a one-standard-deviation change in the axis, in
percentage points.

Counting differs by table. A filing may claim several Nice classes at
once, so a \emph{class-record} counts it once per class while a
\emph{registration} counts it once by serial number; a \emph{debut
filing} is an owner's first, one per owner. A \emph{cohort} is the set of
registrations issued in a calendar year; where a split is instead by
filing year, the text says so.

Two windows bound the samples, each set by an institutional fact. Scoring
runs 1995--2019: a full five-year past reference does not exist before
1990 and a full forward reference does not exist after 2020, and the 1995
cut is about two years more conservative than the measured burn-in
requires (Appendix~\ref{app:data}). Five-year-proof cohorts run
2002--2018: cancellations post at about registration age six and a half,
the record ends in April 2026, so 2018 is the last cohort whose window
has essentially elapsed, and 2002 is the conventional lower cut for this
corpus --- a bound that matters, because carried back to the first
scoreable cohorts the penalty was larger than it has ever been since
(Section~\ref{sec:eras}).

% Neighbours/nomenclature moved to Section 2 (rp_related) for the RP build.
